
\subsection{Optimization model}
\label{section_optimization}

In a REC, there are some usual common parameters. Each member of the community has its own consumption and production profile. Furthermore, some members can have a battery storage system. And as we often want to optimize the self-consumption, the members have a certain number of flexible loads with different kind of flexibility. 

For members that are small consumers at the low voltage level, we can identify different types of standard flexible loads.In \cite{article_6}, the considered flexibilities are of three types : 
\begin{itemize}
    \item{White goods (washing machine, dryer, dishwasher...) with a fixed time of use but flexible starting time}
    \item {Electric vehicle with a charging power that can be controlled but with constraint on the state of charge at departure time}
    \item {Thermal loads that are fully flexible loads}
\end{itemize}
In \cite{article_6}, the way the different flexibilities are managed is associated with a discomfort price in the objective function.

\cite{article_9} consider two main categories of flexible loads : shiftable loads and sheddable loads. Shiftable load are the same as white goods while sheddable loads can be reduced in power for a certain amount of time. So in the previous example, thermal loads and electric vehicles can be considered as sheddable loads.

Finally, in \cite{these_sorel}, they modelize loads in a house for four categories to construct optimization models : 
\begin{itemize}
    \item {Non-flexible loads, they profile is fixed and cannot be changed}
    \item {Shiftable loads, as mentioned before}
    \item {All or nothing loads, that can be turned on or off at any time}
    \item {Electric vehicle loads, here they are considered "All or nothing" when they are at home}
\end{itemize}
By considering all of these flexible loads model it is possible to build optimization models that take into account all the different flexibilities of the members of the REC.

From here on, several parameters can be optimized depending on the goals of the REC. Usually, the cost is minimized as in \cite{article_1, article_7, these_celik}, a discomfort cost as in \cite{article_6} can also be added to the objective function. In \cite{article_5}, the goal is to maximize the robustness of a community that would be disconnected from the main grid. Finally, the ecological impact can also be minimized as in \cite{these_sorel} by minimizing the CO2 emissions associated with the electricity consumption of the community. In \cite{article_7}, they propose a model that take into account demand response requests from the distribution system operator (DSO) in order to provide services to the grid.

In the case where the REC is disconnected from the main grid, what is studied is the operability of the community regarding the uncertainties of production and consumption. \cite{article_3, article_5} both uses robust optimization with a min-max-min approach to ensure the operability of the community in the worst-case scenario.

Several approaches can be used in order to optimize the parameters of the community. In \cite{article_6, article_1} a centralized approach is considered. This means that a central entity has access to all the data of the community and can optimize the parameters for the whole community. Of course, this will reach the best case scenario as all the information is available. However, this can lead to privacy issues and higher computation time. 

To avoid this, \cite{these_celik} proposes different approaches with different levels of data communication. The first approach is the fully centralized one. Then they propose a method where the aggregator communicates the prices to the agents, then the agents propose a first optimal strategy to communicate their surplus and deficits to the aggregator. And then  a loop is initiated between the agents and the aggregator until convergence in order to optimize the battery usage. Finally, a fully decentralized, named "selfish", approach is proposed.


\subsection{Allocation methods within the community}
\label{section_allocation}

With the optimization process, the REC can achieve an optimal planning strategy for exchanging energy within the community and with the main grid. This strategy will lead to overall financial gains for the community. 

In order to distribute these gains, we can use several methods : pricing mechanisms in the community, allocation functions that directly distribute the gains and game-theoretic approaches. Before exploring these methods, the tariff for a single prosumer will be presented.

In this section, the community is made up of a set of agent that represents the community members. Here, each agent is a prosumer, it is associated with a production and a consumption profile that are optimized in the operation phase. A member that does have energy sources will have a null production profile.

% explicité le prix inter communauté 
%Ajouter une subsection si pas de communauté 

\subsubsection{Single prosumer}

Depending on the contract of the prosumer, they will pay a price $p_{buy}(t)$ for each kilowatt-hour consumed. This price might depend on the time of the day. The same way a price for selling is defined $p_{sell}(t)$. Usually, there will also be a subscription fee that depends on the maximum power the prosumer is allowed to use. The subscription fee will not be discussed here as the community will mainly act on the energy prices.

A total cost can then be defined considering the power consumed $P_{cons}$ and produced $P_{prod}$.
\begin{equation}
    Cost = \sum_{t=T_0}^{T_{end}} p_{buy}(t)P_{cons}(t) - p_{sell}(t)P_{prod}(t)
\end{equation}

\subsubsection{Pricing mechanisms}

A common way for allocating the gains is to use pricing mechanisms that will determine the price of energy exchanged within the community. In this case, a local market is created, the prosumers will then buy/sell energy from/to the local market.
\\

\noindent\textbf{Mid-market rate (MMR)}

This method uses the feed-in tariff, which guarantees energy producers a fixed price for renewable electricity fed into the grid under long-term contracts. An important property is to have a fixed reference for the selling price in this method.

The MMR pricing strategy \cite{article_1} is an average price between
the time-of-use and the feed-in tariff. So at each instant we distinguish 3 cases. 


Production of agent $i$ is greater than its consumption, and the total community balance is positive, then : 
% explicité community balance
% spécifier prix en €/kWh
\begin{equation}
    price_{i} = p_{mid} = \frac{p_{buy}+p_{sell}}{2} 
\end{equation} 

Production of agent $i$ is less than its consumption, and the total community balance is positive, then : 

\begin{equation}
    price_{i} = \frac{p_{mid} \cdot \sum\limits_{j \in N}{P_{prod_j}} + p_{buy} \cdot |\sum\limits_{j \in N}{P_{cons_j} - P_{prod_j}|}}{\sum\limits_{j \in N}{P_{cons_j}}} 
\end{equation}
Where $P_{prod_i}$ and $P_{cons_i}$ are respectively the production and consumption of agent $i$. And $N$ is the set of all agents in the community.

Production of agent $i$ is less than its consumption, and the total community balance is negative, then :
\begin{equation}
    price_{i} = \frac{p_{mid} \cdot \sum\limits_{j \in N}{P_{prod_j}} + p_{sell} \cdot |\sum\limits_{j \in N}{P_{cons_j} - P_{prod_j}}|}{\sum\limits_{j \in N}{P_{cons_j}}} 
\end{equation}
\\

\noindent\textbf{Supply-demand rate (SDR)}

The SDR pricing strategy \cite{article_1} sets the internal community price according to a ratio between the community net generation and net demand. At each instant, we have :

\begin{eqnarray}
    rate =& \frac{\sum\limits_{i \in N}{P_{prod_i}}}{\sum\limits_{i \in N}{P_{cons_i}}}  \\
    price =& \begin{cases}
        p_{sell} & \text{if } rate > 1 \\
        p_{buy} & \text{if } rate=0 \\
        p_{sell} rate + (1- rate) & \text{if } 0 < rate \leq 1    
    \end{cases}  
\end{eqnarray}
The price here, is defined as being the same for everyone in the community.
\\

\noindent\textbf{Peer-to-peer transactions}

% Il va ya vaoir des transactions à partir de maintenant donc spécifié le fait que ça change 
% Regrouper dans des chapeux différents qui font sens ce sera plus clair 

% Ajouter la notion d'agent avec consommateur et producteur

In \cite{these_new}, they propose a peer-to-peer pricing mechanism where agents can build bilateral contracts with each other to trade energy. The price of energy is determined by a market mechanism where agents can propose bids and offers. In this study, the agents are paired randomly, but other pairing methods can be used.

\subsubsection{Total profit allocation functions}

Another way of allocating the gains within the community is to only distribute the total profit of the community without using pricing mechanisms. Different methods can be used to do so \cite{these_new, article_1}.
\\

\noindent\textbf{Equal Allocation of Non-Separable Values (EANSV)}

This method \cite{these_new} is based on comparing the individual initial cost without being part of the community and the collective cost of the community. For this it determines the non-separable values that are the values could not be attributed to a single agent. So this method attribute as a cost to each agent its separable cost plus an equal share of the non-separable cost. It can be expressed as follows for an agent $i$ :
\begin{equation}
    C_i^{EANSV} = C_i^{ind} + \frac{C_{community} - \sum\limits_{j \in N}{C_j^{ind}}}{N} 
\end{equation}
Where $C_i^{ind}$ is the individual cost of agent $i$ without being part of the community, $C_{community}$ is the total cost of the community and $N$ is the number of agents in the community.\\

\noindent\textbf{Proportional Allocation (PA)}

This method consist of distributing an equal proportional saving to each agent. It is based on the idea of Nash equilibrium and can be expressed as follows for an agent $i$ :
\begin{equation}
    C_i^{PA} = C_i^{ind} + \frac{|C_i^{ind}|}{\sum\limits_{j \in N}{|C_j^{ind}|}} (C_{community} - \sum\limits_{j \in N}{C_j^{ind}}) \label{eq:PA}
\end{equation}

\subsubsection{Game-theoretic approaches}

This allocation problem can also be studied as a cooperative game. We can then use game theory for computing the distribution. In \cite{article_1,these_sorel}, they propose 2 other methods based on game-theoretic approaches : the Shapley value and the Nucleolus. 

Those two methods uses the principle of coalitions between agents. A coalition $S$ is a subset of $N$, the set of all agents. The value of a coalition $S$ is given by a characteristic function $v(S)$ that gives the total profit of the coalition $S$. \\

\noindent\textbf{Shapley value}

The Shapley value \cite{article_1,these_sorel} is defined as follows for an agent $i$ :
% Rendre les choses plus clair. avec plus d'explication. je peux faire ça avec un exemple.
\begin{equation}
    \phi_i (v) = \frac{1}{n}\sum_{S \subseteq N \setminus \{i\}}{\binom{n-1}{|S|}^{-1} [v(S \cup \{i\}) - v(S)]} \label{eq:shapley}
\end{equation}
with n being the number of element in $N$, $|S|$ is the number of element in set $S$, and $\binom{n-1}{|S|}$ is a binomial coefficient.

This formula can be interpreted as the mean values of the contribution of agent $i$ in each possible coalition without agent $i$.

This method is considered to be fair in a meritocratic sense (see \ref{section_fairness}) as it takes into account the marginal contribution of each agent to all possible coalitions. However, it can be quite complex to compute as the computation time grows exponentially with the number of agents.\\

\noindent\textbf{Nucleolus}

The nucleolus \cite{article_1,these_sorel} of a cooperative game is the solution that maximizes the smallest excess of a coalition. The excess of a coalition S for a given allocation x is defined as :
\begin{equation}
    e(S,x) = v(S) - \sum_{i \in S}{x_i} 
\end{equation}
So the nucleolus can be obtained by solving a serie of linear programming problems that minimize the maximum excess of all coalitions. It is also often simplified using only the prenucleolus, as the nucleolus is quite heavy to compute.

\subsubsection{Other methods}

Other methods have been proposed such as optimization based allocation. Those can be used to propose strategies that answer specific needs of the community. For example, in \cite{these_new}, they propose an optimization framework that aims to take into account the individual preferences of the agents regarding their engagement in the community for ecological purposes.

Dynamic pricing mechanisms using bids and offers can also be used to determine the price of energy exchanged within the community. In this case, several bids strategies can be used.

\subsection{Fairness indexes and equitable distribution}
\label{section-index}

The shapley value and the nucleolus are two methods that aim to provide a fair distribution of the gains among the members of the community. However, there are several ways to define fairness. In \cite{article_1, these_new, article_8, article_2}, they identify different fairness index to evaluate the fairness of a distribution method. % gain, sharing à la place de distribution parce que fonction de distribution

In the following formulas, $x$ is the vector of gain allocations to each agent, and $N$ is the number of agents in the community.

% Specifier d'ou vienne les ref de chaque index

\subsubsection{Jain's index}
Jain's index \cite{article_1, these_new, article_8, article_2} is one of the most common fairness index. It measures the equality of the distribution and is defined as follows :
\begin{equation}
    J(x) = \frac{(\sum_{i=1}^{N}{x_i})^2}{N \cdot \sum_{i=1}^{N}{x_i^2}} 
\end{equation}
Where x is the vector of gain allocations to each agent. The index ranges from 1/N (worst case) to 1 (best case).

\subsubsection{Participation index}
The participation index \cite{article_8} measures the proportion of agents that benefit from the community. It is defined as follows :
\begin{equation}
    PI(x) = \frac{N_{CostReduced}}{N} 
\end{equation}
Where $N_{CostReduced}$ is the number of agents that have a lower cost being part of the community than being alone.

This index can be used to make sure that the distribution method does not disadvantage too many participants.
\subsubsection{MinMax index}
The MinMax index \cite{these_new} measures the ratio between the maximum and minimum allocation among the agents. It aims to evaluate the equality of the distribution in the same way as Jain's index. It is defined as follows :
\begin{equation}
    MM(x) = \frac{min_i{x_i}}{max_i{x_i}} 
\end{equation}

\subsubsection{Quality of Experience (QoE) index}
This index compare the standard deviation to its maximum possible value \cite{these_new}. So the aim is still the same as for the Jain's and MinMax index. It is defined as follows :
\begin{equation}
    QoE(x) = 1 - \frac{\sigma(x)}{max(x_i) - min(x_i)} 
\end{equation}

\subsubsection{Comparison to Shapley value}
The last index measure the difference between the obtained allocation and the Shapley value allocation \cite{article_1, these_new, article_8, article_2}. It is defined as follows :
\begin{equation}
    \Delta = 1 - \sum_{i=1}^{N}{\frac{x_i}{\sum_{k=1}^{N}x_k} - \frac{x_i^{SV}}{\sum_{k=1}^{N}{x_k^{SV}}}} 
\end{equation}
Where $x^{SV}$ is the allocation given by the Shapley value method. The Shapley value is deemed as fair as it calculates the marginal contribution of each agent to all possible coalitions.
% Ajouter 2 trois exemples de ce que je viens de dire.

Using those different fairness indexes, it is possible to evaluate the fairness of a distribution method within a community. This can help choose the most suitable method for a given community depending on its priorities.

\subsection{Social acceptance}
% revenir sur pas mal de chose qu'on a dit 
\cite{article_8} presents a review of the literature on the fairness in local energy systems. They show that fairness withing a community is a key factor in order to put in place systems such as RECs. Furthermore, the concept of fairness is quite subjective and vary depending on the studies that want to take it into account.

They identify 3 main dimensions of fairness : 
\begin{itemize}
    \item {Equality, each participant should receive the same amount of benefits}
    \item {Min-max fairness, the allocation should maximize the minimum benefit received by a participant}
    \item {Meritocracy, the allocation should be proportional to the contribution of each participant}
\end{itemize}

Looking back at the different distribution methods presented previously, the rule based methods such as pricing mechanisms focus on meritocracy. While the redistribution methods can choose more to focus on the aspect they want to prioritize. Then the different indexes seen in \ref{section-index} can be used in order to check the fairness of the final distribution. The same way we can associate the different indexes presented above \ref{section-index} to the different dimensions of fairness. For example, the Jain's index and the MinMax index are more focused on equality, while the comparison to Shapley value is more focused on meritocracy.

\subsection{Sociological profiles of the members of the community}
\label{section_sociological_profiles}

One of the objective of RECs as defined by the European Union is to involve the community in their energy management and in the energy transition. In a more concrete manner, this engagement is meant to provide the needed flexibilities for the community as well as the investment such as PV panels and battery storage. Thus, community engagement is a very important aspect to consider in order for RECs to be successful. \cite{these_sorel} describes 3 main sociological profiles regarding the engagement in a REC :
\begin{itemize}
    \item {The "environmentalist" profile, they are motivated by ecological concerns}
    \item {The "economic" profile, they are motivated by financial gains}
    \item {The "independent" profile, they are motivated by the idea of being autonomous regarding energy}
\end{itemize}
Each member of the community can share more or less one of these profiles. In order to engage the community, it is important to take into account those different profiles while operating the REC. For example, some members might be more willing to accept a higher cost if it leads to a lower ecological impact.

Taking this into account, the optimization process can be adapted using different metrics as objectives. In \cite{these_sorel}, they consider for example self-consumption rate, CO2 emissions and cost as metrics to optimize.